\documentclass[10pt,a4paper]{amsart}
\usepackage{amsmath,amssymb,amsthm,graphicx}
\usepackage{verbatim}
\newenvironment{reference}{\comment}{\endcomment}
%\newenvironment{reference}{}{}
\newenvironment{slogan}{\comment}{\endcomment}
\newenvironment{history}{\comment}{\endcomment}

% For commutative diagrams we use Xy-pic
\usepackage[all]{xy}

% We use 2cell for 2-commutative diagrams.
\xyoption{2cell}
% Theorem environments.
%
\theoremstyle{plain}
\newtheorem{theorem}[subsection]{Theorem}
\newtheorem{proposition}[subsection]{Proposition}
\newtheorem{lemma}[subsection]{Lemma}

\theoremstyle{definition}
\newtheorem{definition}[subsection]{Definition}
\newtheorem{example}[subsection]{Example}
\newtheorem{exercise}[subsection]{Exercise}
\newtheorem{situation}[subsection]{Situation}

\theoremstyle{remark}
\newtheorem{remark}[subsection]{Remark}
\newtheorem{remarks}[subsection]{Remarks}

\numberwithin{equation}{subsection}

% Macros
%
\def\lim{\mathop{\mathrm{lim}}\nolimits}
\def\colim{\mathop{\mathrm{colim}}\nolimits}
\def\Spec{\mathop{\mathrm{Spec}}}
\def\Hom{\mathop{\mathrm{Hom}}\nolimits}
\def\Ext{\mathop{\mathrm{Ext}}\nolimits}
\def\SheafHom{\mathop{\mathcal{H}\!\mathit{om}}\nolimits}
\def\SheafExt{\mathop{\mathcal{E}\!\mathit{xt}}\nolimits}
\def\Sch{\mathit{Sch}}
\def\Mor{\mathop{\mathrm{Mor}}\nolimits}
\def\Ob{\mathop{\mathrm{Ob}}\nolimits}
\def\Sh{\mathop{\mathit{Sh}}\nolimits}
\def\NL{\mathop{N\!L}\nolimits}
\def\CH{\mathop{\mathrm{CH}}\nolimits}
\def\proetale{{pro\text{-}\acute{e}tale}}
\def\etale{{\acute{e}tale}}
\def\QCoh{\mathit{QCoh}}
\def\Ker{\mathop{\mathrm{Ker}}}
\def\Im{\mathop{\mathrm{Im}}}
\def\Coker{\mathop{\mathrm{Coker}}}
\def\Coim{\mathop{\mathrm{Coim}}}

% Boxtimes
%
\DeclareMathSymbol{\boxtimes}{\mathbin}{AMSa}{"02}

%
% Macros for moduli stacks/spaces
%
\def\QCohstack{\mathcal{QC}\!\mathit{oh}}
\def\Cohstack{\mathcal{C}\!\mathit{oh}}
\def\Spacesstack{\mathcal{S}\!\mathit{paces}}
\def\Quotfunctor{\mathrm{Quot}}
\def\Hilbfunctor{\mathrm{Hilb}}
\def\Curvesstack{\mathcal{C}\!\mathit{urves}}
\def\Polarizedstack{\mathcal{P}\!\mathit{olarized}}
\def\Complexesstack{\mathcal{C}\!\mathit{omplexes}}
% \Pic is the operator that assigns to X its picard group, usage \Pic(X)
% \Picardstack_{X/B} denotes the Picard stack of X over B
% \Picardfunctor_{X/B} denotes the Picard functor of X over B
\def\Pic{\mathop{\mathrm{Pic}}\nolimits}
\def\Picardstack{\mathcal{P}\!\mathit{ic}}
\def\Picardfunctor{\mathrm{Pic}}
\def\Deformationcategory{\mathcal{D}\!\mathit{ef}}
\usepackage{cleveref}
\crefname{originallemma}{original lemma}{original lemmas}
\Crefname{originallemma}{Original lemma}{Original lemmas}
\makeatletter
\newlabel{lemma-divided-power-first-order-thickening}{{07HH}{source}{}{sourcepointer-1}{}}
\newlabel{lemma-divided-power-first-order-thickening@cref}{{[originallemma][0][]07HH}{[1][0][]source}{}{}{}}
\makeatother
\title{Divided power differentials via the kernel of the coproduct}
\author{The Stacks Project Authors}
\date{stacks-project:crystalline.tex:local-snapshot-20261012; GNU FDL 1.2 or later}
\begin{document}
\maketitle
\noindent\textit{Editorial note.} All delivered mathematical text is the author's own, copied verbatim from the pinned Stacks Project chapter file crystalline.tex; no statement or proof step was rewritten, reordered or removed.\par\medskip
\noindent\textit{Editorial note.} Span recomputation by label pairing (LANE\_BRIEF 10.2): in this chapter the author's proof environment follows the statement environment as a standalone block; the spec covers the complete environment pair (statement span, proof span) as recorded in the pool candidate.\par\medskip
\ifdefined\hypertarget\hypertarget{sourcepointer-1}{}\fi
\noindent\textit{Source reference.} Original lemma 07HH: sources/stacks-chow/tags.txt, source lines 9552--9552. Cross-chapter standard prerequisite cited by the author inside the counted span; it is used as a pointer and its own proof is not part of this item.\par\medskip
\begin{lemma}
\label{lemma-diagonal-and-differentials}
Let $(A, I, \gamma) \to (B, J, \delta)$ be a homomorphism
of divided power rings. Let $(B(1), J(1), \delta(1))$ be the coproduct
of $(B, J, \delta)$ with itself over $(A, I, \gamma)$, i.e.,
such that
$$
\xymatrix{
(B, J, \delta) \ar[r] & (B(1), J(1), \delta(1)) \\
(A, I, \gamma) \ar[r] \ar[u] & (B, J, \delta) \ar[u]
}
$$
is cocartesian. Denote $K = \Ker(B(1) \to B)$.
Then $K \cap J(1) \subset J(1)$ is preserved by the divided power
structure and
$$
\Omega_{B/A, \delta} = K/ \left(K^2 + (K \cap J(1))^{[2]}\right)
$$
canonically.
\end{lemma}

\begin{proof}
The fact that $K \cap J(1) \subset J(1)$ is preserved by the divided power
structure follows from the fact that $B(1) \to B$ is a homomorphism of
divided power rings.

\medskip\noindent
Recall that $K/K^2$ has a canonical $B$-module structure.
Denote $s_0, s_1 : B \to B(1)$ the two coprojections and consider
the map $\text{d} : B \to K/(K^2 +(K \cap J(1))^{[2]})$ given by
$b \mapsto s_1(b) - s_0(b)$. It is clear that $\text{d}$ is additive,
annihilates $A$, and satisfies the Leibniz rule.
We claim that $\text{d}$ is a divided power $A$-derivation.
Let $x \in J$. Set $y = s_1(x)$ and $z = s_0(x)$.
We will use $\delta$ instead of $\delta(1)$ for
the divided power structure on $J(1)$.
We have to show that $\delta_n(y) - \delta_n(z) = \delta_{n - 1}(y)(y - z)$
modulo $K^2 +(K \cap J(1))^{[2]}$ for $n \geq 1$.
The equality holds for $n = 1$. Assume $n > 1$.
Note that $\delta_i(y - z)$ lies in $(K \cap J(1))^{[2]}$ for $i > 1$.
Calculating modulo $K^2 + (K \cap J(1))^{[2]}$ we have
$$
\delta_n(z) = \delta_n(z - y + y) =
\sum\nolimits_{i = 0}^n \delta_i(z - y)\delta_{n - i}(y) =
\delta_{n - 1}(y) \delta_1(z - y) + \delta_n(y)
$$
This proves the desired equality.

\medskip\noindent
Let $M$ be a $B$-module. Let $\theta : B \to M$ be a divided power
$A$-derivation.
Set $D = B \oplus M$ where $M$ is an ideal of square zero. Define a
divided power structure on $J \oplus M \subset D$ by setting
$\delta_n(x + m) = \delta_n(x) + \delta_{n - 1}(x)m$ for $n > 1$, see
Lemma \ref{lemma-divided-power-first-order-thickening}.
There are two divided power algebra homomorphisms $B \to D$: the first
is given by the inclusion and the second by the map $b \mapsto b + \theta(b)$.
Hence we get a canonical homomorphism $B(1) \to D$ of divided power
algebras over $(A, I, \gamma)$. This induces a map $K \to M$
which annihilates $K^2$ (as $M$ is an ideal of square zero) and
$(K \cap J(1))^{[2]}$ as $M^{[2]} = 0$. The composition
$B \to K/K^2 + (K \cap J(1))^{[2]} \to M$ equals $\theta$ by construction.
It follows that $\text{d}$
is a universal divided power $A$-derivation and we win.
\end{proof}

\end{document}
